\documentclass[11pt,a4paper]{article}
\usepackage[utf8]{inputenc}
\usepackage[spanish]{babel}
\usepackage{amsmath, amssymb}
\usepackage{geometry}
\geometry{top=2.5cm, bottom=2.5cm, left=2.5cm, right=2.5cm}

\begin{document}

\section*{Serie de ejercicios seleccionados para preparar examen de los capítulos 4, 5 y 6}

\begin{itemize}
    \item[\textbf{1.}] Las coordenadas de un objeto que se mueve en el plano $xy$ varían con el tiempo de acuerdo con $x = -(5.00) \operatorname{sen}(\omega t)$ y $y = (4.00) - (5.00) \cos(\omega t)$, donde $\omega$ es una constante, $x$ y $y$ están en metros y $t$ está en segundos. (a) Determine las componentes de velocidad del objeto en $t = 0$. (b) Determine las componentes de la aceleración del objeto en $t = 0$. (c) Escriba expresiones para el vector de posición, el vector velocidad y el vector aceleración del objeto en cualquier tiempo $t > 0$. (d) Describa la trayectoria del objeto en una gráfica $xy$.
    \vspace*{7cm}

    \item[\textbf{2.}] En un bar local, un cliente desliza sobre la barra un tarro de cerveza vacío para que lo vuelvan a llenar. La altura de la barra es $h$. El cantinero está momentáneamente distraído y no ve el tarro, que se desliza de la barra y golpea el suelo a $d$ de la base de la barra. (a) ¿Con qué velocidad el tarro dejó la barra? (b) ¿Cuál fue la dirección de la velocidad del tarro justo antes de golpear el suelo?
    \vspace*{6cm}

    \item[\textbf{3.}] Un proyectil se dispara en tal forma que su alcance horizontal es igual a tres veces su altura máxima. ¿Cuál es el ángulo de proyección?
    \vspace*{6cm}

    \item[\textbf{4.}] La distancia récord en el deporte de lanzar discos es de $81.1 \text{ m}$. Este lanzamiento de discos fue establecido por Steve Urner de Estados Unidos en 1981. Suponiendo que el ángulo de lanzamiento inicial fue de $45^\circ$ y despreciando la resistencia del aire, determine (a) la rapidez inicial del proyectil y (b) el intervalo de tiempo total que voló el proyectil. (c) ¿Cómo cambiaría las respuestas si el alcance fuera el mismo pero el ángulo de lanzamiento fuera mayor que $45^\circ$? Explique.
    \vspace*{6cm}
\end{itemize}

\newpage

\begin{itemize}
    \item[\textbf{5.}] Un proyectil es disparado desde lo alto de un acantilado de altura $h$ sobre el océano. El proyectil se disparó a un ángulo $\theta$ por encima de la horizontal y con una velocidad inicial de $v_i$. (a) Encuentre una expresión simbólica en términos de las variables $v_i, g$ y $\theta$ para el tiempo en que el proyectil alcanza la altura máxima. (b) Utilizando el resultado del inciso (a), encuentre una expresión para la altura máxima $h_{\text{máx}}$ sobre el océano alcanzada por el proyectil en términos de $h, v_i, g$ y $\theta$.
    \vspace*{6cm}

    \item[\textbf{6.}] Un atleta batea una pelota, conectada al extremo de una cadena, en un círculo horizontal. El atleta es capaz de girar la bola con una rapidez de $8.00 \text{ rev/s}$ cuando la longitud de la cadena es $0.600 \text{ m}$. Cuando aumenta la longitud a $0.900 \text{ m}$, puede rotar la pelota solo $6.00 \text{ rev/s}$. (a) ¿Cuál es la tasa de rotación que da la mayor rapidez a la pelota? (b) ¿Cuál es la aceleración centrípeta de la pelota a $8.00 \text{ rev/s}$? (c) ¿Cuál es la aceleración centrípeta en $6.00 \text{ rev/s}$?
    \vspace*{6cm}

    \item[\textbf{7.}] Un neumático de $0.500 \text{ m}$ de radio gira a una rapidez constante de $200 \text{ rev/min}$. Encuentre la rapidez y la aceleración de una pequeña piedra atascada en la banda de rodadura de los neumáticos (en su borde exterior).
    \vspace*{6cm}

    \item[\textbf{8.}] El piloto de un avión registra que la brújula indica un rumbo hacia el oeste. La rapidez del avión respecto al aire es de $150 \text{ km/h}$. El aire se está moviendo con un viento de $30.0 \text{ km/h}$ hacia el norte. Encuentre la velocidad del avión en relación con el suelo.
    \vspace*{6cm}
\end{itemize}

\newpage

\begin{itemize}
    \item[\textbf{9.}] Un automóvil viaja hacia el este con una rapidez de $50.0 \text{ km/h}$. Gotas de lluvia caen con una rapidez constante en vertical respecto de la Tierra. Las trazas de la lluvia en las ventanas laterales del automóvil forman un ángulo de $60.0^\circ$ con la vertical. Encuentre la velocidad de la lluvia en relación con (a) el automóvil y (b) la Tierra.
    \vspace*{6cm}

    \item[\textbf{10.}] Una bola en el extremo de una cuerda gira alrededor en un círculo horizontal de radio $0.300 \text{ m}$. El plano del círculo está a $1.20 \text{ m}$ arriba del suelo. La cuerda se rompe y la bola aterriza a $2.00 \text{ m}$ (horizontalmente) alejándose del punto sobre el suelo directamente debajo de la ubicación de la bola cuando se rompe la cuerda. Encuentre la aceleración radial de la pelota durante su movimiento circular.
    \vspace*{6cm}

    \item[\textbf{11.}] Una partícula parte del origen con velocidad $5\hat{i} \text{ m/s}$ en $t = 0$ y se mueve en el plano $xy$ con una aceleración variable dada por $\vec{a} = (6\sqrt{t}\hat{j})$, donde $\vec{a}$ está en metros por segundo cuadrado y $t$ está en s. (a) Determine el vector velocidad de la partícula como función del tiempo. (b) Determine la posición de la partícula como función del tiempo.
    \vspace*{6cm}

    \item[\textbf{12.}] Un chico lanza una piedra horizontalmente desde la cima de un acantilado de altura $h$ hacia el océano. La piedra golpea el océano a una distancia $d$ desde la base del acantilado. En términos de $h, d$ y $g$, encuentre expresiones para el tiempo $t$ en que la piedra cae en el océano, (b) la rapidez inicial de la piedra, (c) la rapidez de la piedra justo antes de que llegue al océano y (d) la dirección de la velocidad de la piedra justo antes de que llegue al océano.
    \vspace*{8cm}
\end{itemize}

\newpage

\begin{itemize}
    \item[\textbf{13.}] Mientras algún metal fundido salpica, una gota vuela hacia el este con velocidad inicial $v_i$ a un ángulo $\theta_i$ sobre la horizontal y otra gota vuela hacia el oeste con la misma rapidez al mismo ángulo sobre la horizontal, como se muestra en la figura. En términos de $v_i$ y $\theta_i$, encuentre la distancia entre las gotas como función del tiempo.
    \vspace*{6cm}

    \item[\textbf{14.}] La rapidez promedio de una molécula de nitrógeno en el aire es aproximadamente $6.70 \times 10^2 \text{ m/s}$ y su masa es $4.68 \times 10^{-26} \text{ kg}$. (a) Si una molécula de nitrógeno tarda $3.00 \times 10^{-13} \text{ s}$ en golpear una pared y rebotar con la misma rapidez pero moviéndose en la dirección opuesta, ¿cuál es la aceleración promedio de la molécula durante este intervalo de tiempo? (b) ¿Qué fuerza promedio ejerce la molécula sobre la pared?
    \vspace*{6cm}

    \item[\textbf{15.}] La fuerza ejercida por el viento sobre las velas de un bote es de $390 \text{ N}$ hacia el norte. El agua ejerce una fuerza de $180 \text{ N}$ hacia el este. Si el bote (incluyendo su tripulación) tiene una masa de $270 \text{ kg}$, ¿cuáles son la magnitud y dirección de su aceleración?
    \vspace*{6cm}

    \item[\textbf{16.}] La fuerza gravitacional ejercida sobre una pelota de béisbol es $-F_g\hat{j}$. Un pitcher lanza la pelota con velocidad $v\hat{i}$ al acelerarla uniformemente en línea horizontal durante un intervalo de tiempo de $\Delta t = t - 0 = t$. (a) Partiendo del reposo, ¿qué distancia recorre la bola antes de ser liberada? (b) ¿Qué fuerza ejerce el pitcher sobre la pelota?
    \vspace*{6cm}
\end{itemize}

\newpage

\begin{itemize}
    \item[\textbf{17.}] Si un hombre pesa $900 \text{ N}$ en la Tierra, ¿cuál sería su peso en Júpiter, en donde la aceleración en caída libre es de $25.9 \text{ m/s}^2$?
    \vspace*{5cm}

    \item[\textbf{18.}] Un tornillo de hierro de $65.0 \text{ g}$ de masa cuelga de una cuerda de $35.7 \text{ cm}$ de largo. El extremo superior de la cuerda está fijo. Sin tocarlo, un imán atrae el tornillo de modo que permanece fijo, desplazado horizontalmente $28.0 \text{ cm}$ a la derecha desde la línea vertical previa de la cuerda. El imán está ubicado a la derecha del tornillo y al mismo nivel vertical que el tornillo en la configuración final. (a) Dibuje un diagrama de cuerpo libre del tornillo. (b) Encuentre la tensión en la cuerda. (c) Encuentre la fuerza magnética sobre el tornillo.
    \vspace*{7cm}

    \item[\textbf{19.}] Se observa que un objeto de masa $m = 1.00 \text{ kg}$ tiene una aceleración $\vec{a}$ de $10.0 \text{ m/s}^2$ en una dirección a $60.0^\circ$ al noreste, véase la figura. La fuerza $\vec{F}_2$ que se ejerce sobre el objeto tiene una magnitud de $5.00 \text{ N}$ y se dirige al norte. Determine la magnitud y dirección de la fuerza horizontal $\vec{F}_1$ que actúa sobre el objeto.
    \vspace*{6cm}

    \item[\textbf{20.}] Un bloque de $5.00 \text{ kg}$ se coloca encima de otro de $10.0 \text{ kg}$. Una fuerza horizontal de $45.0 \text{ N}$ es aplicada al bloque de $10.0 \text{ kg}$, y el bloque de $5.0 \text{ kg}$ está amarrado a la pared. El coeficiente de fricción cinética entre todos las superficies en movimiento es $0.200$. (a) Dibuje un diagrama de cuerpo libre para cada bloque e identifique las fuerzas de acción-reacción entre los bloques. (b) Determine la tensión en la cuerda y la magnitud de la aceleración del bloque de $10.0 \text{ kg}$.
    \vspace*{7cm}
\end{itemize}

\newpage

\begin{itemize}
    \item[\textbf{21.}] Una curva en una carretera forma parte de un círculo horizontal. Cuando la rapidez de un automóvil que circula por ella es de $14 \text{ m/s}$ constante, la fuerza total horizontal sobre el conductor tiene $130 \text{ N}$ de magnitud. ¿Cuál es la fuerza horizontal total sobre el conductor si la rapidez es $18.0 \text{ m/s}$?
    \vspace*{5cm}

    \item[\textbf{22.}] Un niño de $40.0 \text{ kg}$ se mece en un columpio sostenido por dos cuerdas, cada una de $3.00 \text{ m}$ de largo. La tensión en cada cuerda en el punto más bajo es $350 \text{ N}$. Encuentre (a) la rapidez del niño en el punto más bajo y (b) la fuerza que ejerce el asiento sobre el niño en el punto más bajo. (Desprecie la masa del asiento.)
    \vspace*{6cm}

    \item[\textbf{23.}] Una estudiante está de pie con su mochila en el piso cerca de ella, en un elevador que acelera continuamente hacia arriba con aceleración $a$. El ancho del elevador es $L$. La estudiante da a su mochila una patada rápida en $t = 0$ y le imparte una rapidez $v$ que la hace deslizar a través del piso del elevador. Al tiempo $t$, la mochila golpea la pared opuesta. Encuentre el coeficiente de fricción cinética $\mu_k$ entre la mochila y el piso del elevador.
    \vspace*{7cm}

    \item[\textbf{24.}] La masa de un automóvil deportivo es $1\,200 \text{ kg}$. La forma del cuerpo es tal que el coeficiente de arrastre aerodinámico es $0.250$ y el área frontal es $2.20 \text{ m}^2$. Si ignora todas las otras fuentes de fricción, calcule la aceleración inicial que tiene el automóvil si ha viajado a $100 \text{ km/h}$ y cuando cambia a neutral y se le permite deslizarse.
    \vspace*{6cm}
\end{itemize}

\newpage

\begin{itemize}
    \item[\textbf{25.}] Un pequeño trozo de espuma de poliestireno, material de empaque, se suelta desde una altura de $2.00 \text{ m}$ sobre el suelo. Hasta que llega a su rapidez terminal, la magnitud de su aceleración se conoce mediante $a = g - Bv$. Después de caer $0.500 \text{ m}$, la espuma de estireno en efecto alcanza su rapidez terminal y después tarda $5.00 \text{ s}$ más en llegar al suelo. (a) ¿Cuál es el valor de la constante $B$? (b) ¿Cuál es la aceleración en $t = 0$? (c) ¿Cuál es la aceleración cuando la rapidez es $0.150 \text{ m/s}$?
    \vspace*{7cm}

    \item[\textbf{26.}] Un objeto de masa $m_1 = 4.00 \text{ kg}$ está atado a un objeto de masa $m_2 = 3.00 \text{ kg}$ con la cuerda 1 de longitud $\ell = 0.500 \text{ m}$. La combinación se hace oscilar en una trayectoria circular vertical sobre una segunda cuerda, la cuerda 2, de longitud $\ell = 0.500 \text{ m}$. Durante el movimiento, las dos cuerdas son colineales en todo momento como se muestra en la figura. En la parte superior de su movimiento, $m_2$ viaja a $v = 4.00 \text{ m/s}$. (a) ¿Cuál es la tensión en la cuerda 1 en este instante? (b) ¿Cuál es la tensión en la cuerda 2 en este instante? (c) ¿Cuál cuerda se romperá primero si la combinación se gira más y más rápido?
    \vspace*{7cm}

    \item[\textbf{27.}] Un camión se mueve con aceleración constante $a$ en una colina que hace un ángulo $\phi$ con la horizontal como en la figura. Una pequeña esfera de masa $m$ está suspendida desde el techo de la camioneta por un cable ligero. Si el péndulo hace un ángulo constante $\theta$ con la perpendicular al techo, ¿a qué es igual $a$?
    \vspace*{6cm}

    \item[\textbf{28.}] Un modelo de avión de masa $0.750 \text{ kg}$ vuela con una rapidez de $35.0 \text{ m/s}$ en un círculo horizontal en el extremo de un cable de control de $60.0 \text{ m}$ de largo. En la figura se muestran las fuerzas ejercidas sobre el avión: la tensión en el alambre de control, la fuerza gravitacional y la elevación aerodinámica que actúa en $\theta = 20.0^\circ$ hacia adentro de la vertical. Calcule la tensión en el cable, suponiendo que hace un ángulo constante de $\theta = 20.0^\circ$ con la horizontal.
    \vspace*{6cm}

    \item[\textbf{29.}] Para $t < 0$, un objeto de masa $m$ no experimenta fuerza y se mueve en la dirección $x$ positiva con una rapidez constante $v_i$. Comenzando en $t = 0$, cuando el objeto pasa de la posición $x = 0$, experimenta una fuerza neta resistiva proporcional al cuadrado de su rapidez: $\vec{F}_{\text{neta}} = -mkv^2 \hat{i}$, donde $k$ es una constante. La rapidez del objeto después de $t = 0$ viene dada por $v = v_i / (1 + kv_it)$. (a) Encuentre la posición $x$ del objeto como una función del tiempo. (b) Encuentre la velocidad del objeto como una función de la posición.
    \vspace*{7cm}
\end{itemize}

\end{document}
